\section{Setting: what a post-search score estimates}
\label{sec:setting}
\textbf{Setup.} A \emph{configuration} $k\in\{1,\dots,K\}$ is a complete prompt or LLM-program specification (instruction, demonstrations, output format, \dots). For a question $x$ from a benchmark distribution $P$, $Y_k(x)\in\{0,1\}$ is the correctness of the system under $k$ (with deterministic decoding; stochastic decoding only adds noise); $\theta_k=\E_P Y_k$ is its \emph{true accuracy}. A developer holds $n$ labelled questions $D$ drawn i.i.d.\ from $P$ and runs a search procedure $\mathcal A$ that queries outcomes $Y_k(x)$, $x\in D$, and returns a configuration $\khat$. The common practice we examine is to report the \emph{naive score} $\hat S=\frac1n\sum_{x\in D}Y_{\khat}(x)$. The target of interest is $\theta_{\khat}$, the accuracy of the deployed configuration on fresh questions; the \emph{gap} is $\hat S-\theta_{\khat}$ and the \emph{optimism} is $\mathcal O=\E[\hat S-\theta_{\khat}]$.

\textbf{Exhaustive search.} Let $\bar Y_k$ be the dev accuracy of $k$, $\epsilon_k=\bar Y_k-\theta_k$, $\khat=\arg\max_k\bar Y_k$, $\theta^\star=\max_k\theta_k$.
\begin{proposition}[Size of the optimism]\label{prop:opt}
(a) $\mathcal O\;\ge\;\theta^\star-\E[\theta_{\khat}]\;\ge 0$: the optimism is at least the expected regret of the search. (b) $\mathcal O\le \E\max_k\epsilon_k\le\sqrt{\ln K/(2n)}$. (c) If all $\theta_k$ are equal and $(\epsilon_k)$ is Gaussian with common variance $\sigma_n^2=\bar p(1-\bar p)/n$ and equal pairwise correlation $\rho$ (the correlation of per-question outcomes between configurations), then $\mathcal O=\sigma_n\sqrt{1-\rho}\;m_K$, with $m_K=\E\max_{k\le K}Z_k\le\sqrt{2\ln K}$ the expected maximum of $K$ i.i.d.\ standard normals ($m_{2}{=}0.56$, $m_{10}{=}1.54$, $m_{40}{=}2.16$, $m_{80}{=}2.43$).
\end{proposition}
\emph{Proof.} (a) $\E\hat S=\E\max_k\bar Y_k\ge\max_k\E\bar Y_k=\theta^\star$. (b) $\hat S-\theta_{\khat}=\epsilon_{\khat}\le\max_k\epsilon_k$; each $\epsilon_k$ is sub-Gaussian with variance proxy $1/(4n)$ (Hoeffding), so the maximal inequality gives $\sqrt{2\cdot\frac1{4n}\ln K}$. (c) Write $\epsilon_k=\sigma_n(\sqrt\rho Z_0+\sqrt{1-\rho}Z_k)$; the common part has mean zero and $\hat S - \theta = \max_k \epsilon_k$ since all $\theta_k$ coincide. $\square$

The optimism therefore (i) shrinks only as $n^{-1/2}$, (ii) grows with the number of configurations \emph{effectively} compared, $\sqrt{(1-\rho)}\,m_K$, so that highly correlated search spaces (as are small variations of one prompt) are less inflated than diverse ones, and (iii) is largest in units of $\sigma_n$ when configurations are equally good and decays as the true spread $\mathrm{sd}(\theta_k)$ grows relative to $\sigma_n$ (a ``signal-to-noise'' ratio $\mathrm{SNR}=\mathrm{sd}(\theta_k)/\sigma_n$): when one configuration is clearly best, selection is stable and there is little to be optimistic about. We check (c) and the SNR dependence in \cref{sec:exp}. For adaptive optimizers the same logic applies with the information the search extracts from $D$, which is more limited for budget-constrained searches (successive halving, minibatch search) than for exhaustive evaluation.

\begin{proposition}[Budget-confounded comparisons]\label{prop:cmp}
Let systems A and B be tuned on the same $D$ over disjoint spaces with $K_A>K_B$ equally good, equally correlated configurations. Then $\E[\hat S_A-\hat S_B]-(\E\theta_{\khat_A}-\E\theta_{\khat_B})=\sigma_n\sqrt{1-\rho}\,(m_{K_A}-m_{K_B})>0$.
\end{proposition}
\noindent The reported advantage of the system that was tuned harder is inflated even if tuning is equally effective for both; \cref{sec:exp} measures this on real outcomes.
